% ĐỀ KIỂM TRA CẤP 3 CÓ TỰ LUẬN: PHẦN I, II, III trắc nghiệm, PHẦN IV tự luận. In đen trắng.
% Cách dùng (b): thư mục soạn bài có mẫu con mỏng de10.sty (= kiem_tra_thpt + \mon{TOÁN 10}).
\documentclass[12pt]{article}
\usepackage{de10}
\begin{document}
\kythi{Kiểm tra giữa học kì I}
\made{101}
\taodaude

\phantn
\begin{tn}
Giá trị của $\sin\dfrac{\pi}{6}$ bằng
\pa{\dung{$\dfrac12$}}{$\dfrac{\sqrt3}{2}$}{$1$}{$0$}
\end{tn}

\phands
\begin{ds}
Cho cấp số cộng $(u_n)$ có $u_1=2$, công sai $d=3$.
\ydung{$u_2=5$.}
\ysai{$u_5=15$.}
\ydung{$S_4=26$.}
\ysai{$(u_n)$ là dãy số giảm.}
\end{ds}

\phantln
\begin{tln}
Giá trị nhỏ nhất của hàm số $y=x^2-4x+5$ bằng bao nhiêu?
\traloi{1}
\end{tln}

\phantl[2,0]
\begin{bt}\diem{1,0}
Giải phương trình $\sin x=\dfrac12$.
\loigiai{$x=\dfrac{\pi}{6}+k2\pi$ hoặc $x=\dfrac{5\pi}{6}+k2\pi$.}
\end{bt}
\begin{bt}\diem{1,0}
Tính tổng $1+3+5+\cdots+99$.
\loigiai{Có $50$ số hạng, tổng $=50^2=2500$.}
\end{bt}
\bangdapan
\hetde
\end{document}
